\documentclass[instructions.tex]{subfiles}
\begin{document}
 
\chapter{SUITE DU MÊME SUJET.}
 
Respect dû aux personnes consacrées à Dieu.
 
\primo Un libertin dit un jour à un homme de bien : De quoi servent au monde tant de prêtres, tant de religieux et de religieuses ? Et vous, lui répondit-il, de quoi y servez-vous ? de quoi y servent tant de mondains et de scandaleux ? de quoi servent tant d'hommes avides, ambitieux, qui accablent les pauvres et le peuple ?
 
De quoi servent, dites-vous, tant de prêtres et de personnes religieuses ? Ils servent, dit saint Augustin, à louer Dieu, tandis que vous l'offensez ; à faire pénitence, tandis que vous vivez dans le crime. Ils servent, disait l'empereur Justinien, à apaiser la colère de Dieu, par leurs prières, sur les empires et les royaumes ; car le monde périrait par les fléaux du ciel, s'il n'y avait pas des âmes saintes sur la terre. Ils servent enfin à justifier la conduite du Tout-Puissant, et à vous condamner.
 
Que direz-vous au jugement de Dieu, à la vue de tant de personnes plus délicates et plus riches que vous, qui ont tout quitté pour vivre dans la pénitence et la retraite ? Que direz-vous, lorsque vous verrez des filles de qualité qui ont sacrifié les plus brillantes fortunes pour aller dans un cloître crucifier leur chair, et pleurer les péchés du monde, tandis que vous passez votre vie dans le désordre ? Que direz-vous à la vue de tant de ministres du Seigneur qui tâchent de sauver les âmes en les retirant du vice et de l'erreur, tandis que vous les perdez ? Le démon ne peut souffrir que Dieu soit honoré ; c'est pour cela qu'il tâche d'inspirer du mépris pour les personnes sacrées et les serviteurs de Dieu.
 
\secundo Les prêtres et les religieux ne sont pas impeccables. Prétendre qu'ils soient sans défauts, c'est vouloir qu'ils ne soient point hommes. Mais si quelques-uns s'écartent de leur devoir, faudrait-il pour cela condamner ceux qui sont innocens ? Parce qu'un officier est infidèle à son prince, faut-il mépriser les autres, et conclure que tous les officiers sont des perfides ? Parce que Judas a trahi Jésus-Christ, faut-il conclure que les autres disciples étaient des traîtres et des apostats ?
 
Pour un prêtre qui oublierait son devoir, combien d'autres qui vivent en saints, en pasteurs désintéressés et pleins de zèle ! Pour un religieux qui s'écarterait de sa règle, combien d'autres religieux qui sont exemplaires et pénitens ! Pourquoi mépriser ceux-ci pour les fautes de ceux-là ?
 
Un prêtre qui vit selon la sainteté de son état, une personne religieuse qui vit selon l'esprit de sa règle, donnent plus de gloire à Dieu, font plus de bonnes œuvres et d'actes de religion dans une semaine, que les gens du monde les mieux réglés n'en font dans un mois.
 
Après tout, si un prêtre ou un religieux manquait à quelques devoirs, faut-il que les gens du monde manquent de charité pour lui, de vénération pour tous les apôtres ; ainsi quand, par un grand malheur, des hommes d'un auguste caractère viennent à tomber, tremblez pour vous ; dites en vous-même. Si l'on fait des chutes dans un état aussi saint, que dois-je craindre au milieu des dangers du monde ?
 
Priez pour eux. Plus ils sont élevés par leur caractère, plus aussi sont-ils exposés aux dangers et dignes de compassion. Cachez leurs défauts, n'en parlez point. Si je voyais un ministre de Dieu tomber dans une faute, disait l'empereur Constantin, loin de la publier, j'irais moi-même le couvrir de mon manteau impérial. Il est de l'intérêt du démon de décrier les personnes consacrées à Dieu ; il n'y réussit que trop par la langue des méchans.
 
Êtes-vous innocent, vous qui parlez des oints du Seigneur ? Il vous sied mal de les noircir, vous qui êtes encore plus criminel. Celui d'entre vous qui est sans péché, dit Jésus-Christ aux pharisiens, qu'il jette la première pierre. La punition ordinaire de ceux qui méprisent les pasteurs, c'est de mourir sans prêtres et sans sacremens. Vous méprisez les ministres de Dieu, Dieu vous méprisera à son tour\footnote{Væ qui spernis, nonne et ipse sperneris? Is. 33. 1.}.
 
\section{Résolutions}
 
\primo Je ne relèverai point les défauts ni les fautes des prêtres et des autres personnes consacrées à Dieu.

\secundo S'il arrive dans le plus saint des états quelque désordre et quelque scandale, j'en gémirai devant Dieu ; je lui en demanderai pardon ; je prendrai garde à moi et je penserai que c'est une épreuve à laquelle Dieu a permis que la foi des fidèles fût soumise. 

\prayer{O mon Dieu ! qui avez voulu que les hommes fussent conduits au salut par d'autres hommes, donnez à nos guides spirituels une sainteté qui brille comme une grande lumière aux yeux de tous ceux dont ils sont chargés.}
 
\end{document}
